\section{The Haar hard edge and the quenched law}
\label{sec:hard-edge}

Let \(\mathsf U_g\) be Haar distributed in \(\USp(2g)\), and write its positive
eigenangles as
\(0<\vartheta_1<\cdots<\vartheta_g<\pi\).  With \(d=2g+1\), set
\[
 \Xi_g=\sum_{j=1}^g\delta_{d\vartheta_j/(2\pi)},
\]
regarded as a point measure on the \emph{closed} half-line
\([0,\infty)\).  Keeping the endpoint is important: vague convergence on
\((0,\infty)\) would not see a point moving to zero.

\subsection{The finite kernel and its closed-edge limit}

Weyl integration, followed by the change of variables
\(y=d\vartheta/(2\pi)\), gives the projection kernel
\begin{equation}
 \widetilde K_g(x,y)=\frac4d\sum_{k=1}^g
   \sin\frac{2\pi kx}{d}\sin\frac{2\pi ky}{d},
 \qquad 0\le x,y\le \frac d2 .
 \label{eq:finite-sp-kernel}
\end{equation}
Indeed, before rescaling the orthonormal system is
\((2/\pi)^{1/2}\sin(k\vartheta)\), \(1\le k\le g\).  Thus no
normalizing constant is hidden in \eqref{eq:finite-sp-kernel}.  For every
fixed \(L\), the right side is a Riemann sum, uniformly on
\([0,L]^2\), and
\begin{equation}
 \widetilde K_g(x,y)\longrightarrow
 \frac2\pi\int_0^\pi\sin(tx)\sin(ty)\,dt
 =K_{\mathrm{Sp}}(x,y).
 \label{eq:kernel-riemann-limit}
\end{equation}
This convergence includes \(x=0\) and \(y=0\).  Moreover,
\begin{equation}
 0\le \widetilde K_g(y,y)\le \frac{4g}{2g+1}<2,
 \qquad
 K_{\mathrm{Sp}}(y,y)
 =1-\frac{\sin(2\pi y)}{2\pi y}<2,
 \label{eq:kernel-diagonal-envelope}
\end{equation}
where the last expression has value zero at \(y=0\).

Extend each finite kernel by zero off \([0,d/2]\).  Local uniform convergence
first gives convergence of the corresponding operators on compactly
supported \(L^2\) test functions.  Because all the finite operators and the
limit are orthogonal projections, hence contractions, density upgrades this
to weak-operator convergence on \(L^2([0,\infty))\).  Local uniform convergence
of the diagonals and the common bound in
\eqref{eq:kernel-diagonal-envelope} also give convergence of every bounded
local trace.  The limiting kernel is a locally trace-class orthogonal
projection, so the determinantal process exists, and Soshnikov's operator
criterion gives~\cite[Theorems~3 and~5]{Soshnikov2000}
\begin{equation}
 \Xi_g\Longrightarrow\PiSp
 \quad\text{vaguely on }[0,\infty).
 \label{eq:closed-edge-process-convergence}
\end{equation}
One can also see tightness directly from
\(\E\Xi_g([0,L])\le2L\).  In particular,
\begin{equation}
 \sup_g\E\Xi_g([0,\eta])\le2\eta,
 \label{eq:no-escape-zero}
\end{equation}
which records explicitly that mass cannot be lost through the hard edge.
The limit has no point at zero, since its intensity there is zero.

\subsection{Critical coefficients and their tails}

For a rescaled angle \(y\in[0,d/2]\), the coefficient in the finite
random-phase law is exactly
\begin{align}
 a_{g,\lambda}(y)
 &=\frac{4(1-e^{-\lambda/d})}
 {\sqrt{(1-e^{-\lambda/d})^2
       +4e^{-\lambda/d}\sin^2(\pi y/d)}} \notag\\
 &=\frac{4\sinh(\lambda/(2d))}
 {\sqrt{\sinh^2(\lambda/(2d))+\sin^2(\pi y/d)}}.
 \label{eq:finite-critical-coefficient}
\end{align}
On every compact subset of the closed half-line,
\begin{equation}
 a_{g,\lambda}(y)\longrightarrow
 a_\lambda(y)=\frac{4\lambda}
 {\sqrt{\lambda^2+(2\pi y)^2}}
 \label{eq:coefficient-local-limit}
\end{equation}
uniformly.  There is also a uniform global envelope.  Since
\(\sin(\pi y/d)\ge2y/d\) for \(0\le y\le d/2\), and \(d\ge5\),
\begin{equation}
 a_{g,\lambda}(y)\le
 \min\left\{4,\frac{A_\lambda}{y}\right\},
 \qquad A_\lambda:=\lambda e^{\lambda/10},
 \label{eq:coefficient-envelope}
\end{equation}
for \(y>0\).  The limiting coefficient obeys the same bound (in fact
\(a_\lambda(y)\le2\lambda/(\pi y)\)).  Combining
\eqref{eq:kernel-diagonal-envelope} and
\eqref{eq:coefficient-envelope}, for every \(R>0\),
\begin{equation}
 \sup_g\E\sum_{y\in\Xi_g,\ y>R}a_{g,\lambda}(y)^2
 \le \frac{2A_\lambda^2}{R},
 \qquad
 \E\sum_{y\in\PiSp,\ y>R}a_\lambda(y)^2
 \le \frac{2A_\lambda^2}{R}.
 \label{eq:coefficient-tail}
\end{equation}
Here and below a sum over a point measure counts multiplicity.  Notice also
that
\[
 \int_0^\infty a_\lambda(y)^2\,dy=4\lambda.
\]
Consequently the limiting squared coefficient sum is finite almost surely,
and Markov's inequality turns the first bound in
\eqref{eq:coefficient-tail} into the uniform probability estimate
\(2A_\lambda^2/(\eta R)\) at threshold \(\eta>0\).

We now apply the closed-hard-edge stability result proved in
\cref{app:marked-laws}.  Extend \(a_{g,\lambda}\) continuously beyond
\(d/2\) in any way; this does not change its values on \(\Xi_g\).  The
local limit \eqref{eq:coefficient-local-limit} and the tail estimate
\eqref{eq:coefficient-tail} are precisely the two coefficient hypotheses of
\cref{prop:marked-stability}.  Coupling the same parity variable and the same
phases gives
\begin{equation}
 \Wtwo^2\bigl(T_{g,\lambda}(\mathsf U_g),\nu_\lambda(\PiSp)\bigr)
 \le \frac12\sum_{\epsilon=0}^1
       |c_{g,\lambda}^{(\epsilon)}-1|^2
     +\frac12\sum_{j\ge1}|u_{g,j}-u_j|^2
 \label{eq:haar-shared-coupling}
\end{equation}
on a Skorokhod coupling, where the two coefficient vectors are matched in
increasing point order and padded by zeros.  The right side tends to zero in
probability.  This proves
\begin{proposition}[Haar quenched limit]
\label{prop:haar-quenched-limit}
For Haar \(\mathsf U_g\in\USp(2g)\),
\[
 T_{g,\lambda}(\mathsf U_g)\Longrightarrow\nu_\lambda(\PiSp)
 \quad\text{in }(\Ptwo,\Wtwo).
\]
The random measure on the right is Borel measurable and almost surely
absolutely continuous.
\end{proposition}

For clarity, the last assertion does not rest only on the formal stability
lemma.  The limiting process is locally finite and has no point at zero.  It
is also infinite almost surely.  If
\(N_L=\PiSp([0,L])\), the determinantal count identities give
\[
 \E N_L=L+O(1),
 \qquad \Var N_L\le \E N_L.
\]
For each fixed integer \(k\), Chebyshev's inequality yields
\(\Pr(N_L\le k)\to0\).  Monotonicity in \(L\) proves infinitude.  The
configuration therefore has a first point \(y_1>0\).  Conditional on the
configuration, the series in \eqref{eq:quenched-limit-law} converges almost
surely and in \(L^2\); its first summand has a nonzero arcsine density and is
independent of the convergent remainder.  Their convolution is absolutely
continuous.

\subsection{Two distinct nondegeneracy statements}

We first prove that the random \emph{probability measure} is not deterministic.
Every conditional law has mean one and conditional variance
\begin{equation}
 V_\lambda(\xi)=\frac12\sum_{y\in\xi}a_\lambda(y)^2
 =\sum_{y\in\xi} f_\lambda(y),
 \qquad
 f_\lambda(y)=\frac{8\lambda^2}{\lambda^2+4\pi^2y^2}.
 \label{eq:conditional-variance-statistic}
\end{equation}
The representation in \eqref{eq:kernel-riemann-limit} identifies
\(K_{\mathrm{Sp}}\) as the orthogonal projection, under the unitary sine
transform, onto frequencies in \([0,\pi]\).  Hence
\begin{equation}
 \int_0^\infty|K_{\mathrm{Sp}}(x,y)|^2\,dy
 =K_{\mathrm{Sp}}(x,x).
 \label{eq:projection-kernel-identity}
\end{equation}
Since \(f_\lambda\in L^1(0,\infty)\cap L^2(0,\infty)\), compactly
supported approximation in the standard determinantal covariance identity
is legitimate~\cite[Definition~3 and (1.20)]{Soshnikov2000}.  Explicitly,
using \(K_{\mathrm{Sp}}(x,x)<2\) and
\(\rho_2(x,y)\le K_{\mathrm{Sp}}(x,x)K_{\mathrm{Sp}}(y,y)<4\),
\[
 \E\left(\sum_{y\in\PiSp,\ y>R}f_\lambda(y)\right)^2
 \le 2\int_R^\infty f_\lambda(y)^2\,dy
    +4\left(\int_R^\infty f_\lambda(y)\,dy\right)^2
 \longrightarrow0.
\]
Passing to the limit from compact truncations therefore gives
\begin{equation}
 \Var V_\lambda(\PiSp)
 =\frac12\int_0^\infty\!\int_0^\infty
   (f_\lambda(x)-f_\lambda(y))^2
   |K_{\mathrm{Sp}}(x,y)|^2\,dx\,dy.
 \label{eq:projection-dpp-variance}
\end{equation}
All factors in this formula are essential, including \(1/2\).  The integral
is finite because it is at most
\(2\int f_\lambda(x)^2K_{\mathrm{Sp}}(x,x)\,dx\).  It is strictly
positive: \(f_\lambda\) is strictly decreasing on \((0,\infty)\), while
the real-analytic kernel is nonzero on an open set off the diagonal.  We have
proved the following.

\begin{proposition}[Nondegeneracy of the quenched measure]
\label{prop:random-law-nondegenerate}
The random element \(\nu_\lambda(\PiSp)\) of \(\Ptwo\) is not almost surely
constant.
\end{proposition}

Indeed, if that conditional law were fixed, then its variance would be
fixed, contrary to \eqref{eq:projection-dpp-variance}.  This argument says
nothing by itself about the single scalar
\(\nu_\lambda(\PiSp)((0,\infty))\); we prove that assertion separately.

\begin{proposition}[Nondegeneracy of the half-line mass]
\label{prop:scalar-nondegenerate}
The random variable
\[
 F_\lambda(\PiSp):=\nu_\lambda(\PiSp)((0,\infty))
\]
is not almost surely constant.
\end{proposition}

\begin{proof}
For a fixed nonempty square-summable configuration, write
\(S_\xi=\sum_{y\in\xi}a_\lambda(y)\cos\Phi_y\).  Its law is symmetric and,
by splitting off the first summand, absolutely continuous.  Moreover
\begin{equation}
 \Pr_\Phi(|S_\xi|<1)>0.
 \label{eq:positive-small-ball}
\end{equation}
To see this, choose a finite head so that the tail variance is small enough
 that \(\Pr(|S_{\mathrm{tail}}|<1/2)>0\), and then require all the finitely many
head cosines to be sufficiently close to zero.  Symmetry and absolute
continuity now give
\begin{equation}
 F_\lambda(\xi)
 =\frac12\left(1+\Pr_\Phi(|S_\xi|<1)\right)>\frac12.
 \label{eq:scalar-strict-half}
\end{equation}

Choose a compact interval \(I\subset(0,\infty)\), bounded away from zero
but close enough to it that
\(a_\lambda(y)\ge a_0>2\) on \(I\).  For every integer \(N\),
\begin{equation}
 \Pr\{\PiSp(I)\ge N\}>0.
 \label{eq:dpp-crowding}
\end{equation}
Indeed, the factorial-moment formula
\cite[Definition~2 and (1.5)]{Soshnikov2000} is
\[
 \E(\PiSp(I))_N
 =\int_{I^N}\det\bigl(K_{\mathrm{Sp}}(x_i,x_j)\bigr)_{i,j\le N}
   \,dx_1\cdots dx_N.
\]
For distinct positive \(x_i\), the determinant is the Gram determinant in
\(L^2(0,\pi)\) of the functions
\(t\mapsto(2/\pi)^{1/2}\sin(tx_i)\).  These functions are linearly independent:
differentiate an identity \(\sum_i c_i\sin(tx_i)=0\) at zero and use the
Vandermonde matrix in \(x_i^2\).  The determinant, hence the factorial moment,
is positive.  This proves \eqref{eq:dpp-crowding}.

On the event in \eqref{eq:dpp-crowding}, select the first \(N\) points of
\(\PiSp\cap I\).  Their coefficients lie in \([a_0,4]\), and their marked
sum has characteristic function
\(\prod_{j=1}^N J_0(a_jt)\).  Uniformly for \(a_j\in[a_0,4]\) and
\(N\ge3\), Fourier inversion gives
\begin{equation}
 \left\|p_{\sum_{j=1}^Na_j\cos\Phi_j}\right\|_\infty
 \le \frac{C_{\lambda,I}}{\sqrt N}.
 \label{eq:uniform-bessel-concentration}
\end{equation}
Here are the constants and the inversion step.  Uniformly for
\(a\in[a_0,4]\), the power series at zero supplies \(c,\delta>0\) such that
\(|J_0(at)|\le e^{-ct^2}\) for \(|t|\le\delta\).  The strict inequality
\(|J_0(x)|<1\) for \(x\ne0\) gives
\[
 r:=\sup_{a\in[a_0,4],\ \delta\le|t|\le T}|J_0(at)|<1
\]
for every fixed \(T>\delta\).  Finally choose an absolute \(B\) and then
\(T\) large enough that the standard large-argument bound
\(|J_0(x)|\le Bx^{-1/2}\) applies for \(x\ge a_0T\) and
\(q:=B/\sqrt{a_0T}<1\).  On the positive half-line,
\begin{align*}
 \int_0^\delta e^{-cNt^2}\,dt
   &\le\frac{\sqrt\pi}{2\sqrt{cN}},\\
 \int_\delta^T r^N\,dt&\le Tr^N,\\
 \int_T^\infty\left(\frac{B}{\sqrt{a_0t}}\right)^Ndt
   &=\frac{T}{N/2-1}q^N \qquad(N>2).
\end{align*}
Doubling these bounds proves that the characteristic function is in
\(L^1(\mathbb R)\), with norm \(O_{\lambda,I}(N^{-1/2})\).  Fourier inversion
therefore gives a continuous density bounded by
\((2\pi)^{-1}\) times that norm, proving
\eqref{eq:uniform-bessel-concentration}.

Convolving with the independent contribution of all other points cannot
increase the density norm.  Hence, on the positive-probability crowding
event,
\begin{equation}
 \frac12<F_\lambda(\PiSp)
 \le\frac12+\frac{C_{\lambda,I}}{\sqrt N}.
 \label{eq:crowding-upper-bound}
\end{equation}
If \(F_\lambda(\PiSp)=c\) almost surely, then
\eqref{eq:scalar-strict-half} gives \(c>1/2\).  Choosing \(N\) so large
that the right side of \eqref{eq:crowding-upper-bound} is smaller than
\(c\) contradicts \eqref{eq:dpp-crowding}.  Thus the scalar is nonconstant.
\end{proof}
